% GENERATED FILE. Edit canonical lessons, then run npm run generate:books.
% canonical-source-hash: fnv1a64:3ac970ea18c2c538
% canonical-lessons: BN-C289-jolochchhvas, BN-C289-toliye-jaoya, BN-C289-uchu, BN-C289-bhela, BN-C289-plabito-hooya

\chapter{The Water Rises}
\label{ch:bn-the-water-rises}
\hlchaptermodality{289}

\begin{chapteropening}
\textbf{By the end of this chapter:} \emph{I can describe a tidal surge and the rising water.}

Complete the last lesson of chapter 289: I can describe a tidal surge and the rising water.
\end{chapteropening}

\section[\texorpdfstring{\bn{জলোচ্ছ্বাস} (jôlochchhvās)}{জলোচ্ছ্বাস (jôlochchhvās)}]{\bn{জলোচ্ছ্বাস} (jôlochchhvās) — a tidal surge}

\label{lesson:BN-C289-jolochchhvas}



\begin{quote}
Before the new one: say the Bengali for to uproot, then the Bengali for electricity.
\end{quote}

\subsection*{You'll want to know: \bn{জলোচ্ছ্বাস}}
\textbf{\bn{জলোচ্ছ্বাস}} — \emph{jôlochchhvās} — \textquotedblleft{}a tidal surge\textquotedblright{}.

\begin{quote}
\bn{ঘূর্ণিঝড়ের} \bn{সঙ্গে} \bn{একটা} \bn{বড়} \bn{জলোচ্ছ্বাস} \bn{এল।}
\end{quote}

\begin{quote}
\emph{ghūrnijhôṛer sôngge ekṭā bôṛo jôlochchhvās elô.}
\end{quote}

\begin{quote}
A great tidal surge came with the cyclone.
\end{quote}

\begin{cousinweb}
\textbf{\bn{জলোচ্ছ্বাস}} (\emph{jôlochchhvās}) joins \textbf{\bn{জল}}, water, to Sanskrit \emph{ucchvāsa}, 'a breathing out, a deep breath' (Monier-Williams): the sea heaving up over the land.
\end{cousinweb}

\begin{grammarlens}[title={What you've built}]
One more word for telling what happened.
\end{grammarlens}

\subsection*{Your turn}
\begin{itemize}
\raggedright
  \item \emph{Say it:} \emph{jôlochchhvās}
  \item \emph{Say it:} \emph{jôlochchhvās}, once more
  \item \emph{Say it:} \emph{ghūrnijhôṛer sôngge ekṭā bôṛo jôlochchhvās elô.} — A great tidal surge came with the cyclone.
  \item \emph{Recall:} say the Bengali for to uproot, then the Bengali for electricity, then say \emph{jôlochchhvās} again
\end{itemize}

\begin{tcolorbox}[breakable,colback=teal!4,colframe=teal!35!black,title={Before you move on}]
What does \bn{জলোচ্ছ্বাস} mean? (\textquotedblleft{}A tidal surge\textquotedblright{}.) Say it once more.
\end{tcolorbox}
\section[\texorpdfstring{\bn{তলিয়ে} \bn{যাওয়া} (tôliye jāoyā)}{তলিয়ে যাওয়া (tôliye jāoyā)}]{\bn{তলিয়ে} \bn{যাওয়া} (tôliye jāoyā) — to go under; to be submerged}

\label{lesson:BN-C289-toliye-jaoya}



\begin{quote}
Before the new one: say the Bengali for electricity, then the Bengali for a tidal surge.
\end{quote}

\subsection*{You'll want to know: \bn{তলিয়ে} \bn{যাওয়া}}
\textbf{\bn{তলিয়ে} \bn{যাওয়া}} — \emph{tôliye jāoyā} — \textquotedblleft{}to go under; to be submerged\textquotedblright{}.

\begin{quote}
\bn{জলোচ্ছ্বাসে} \bn{নিচু} \bn{গ্রামগুলো} \bn{তলিয়ে} \bn{গেল।}
\end{quote}

\begin{quote}
\emph{jôlochchhvāse nichu grāmgulo tôliye gelô.}
\end{quote}

\begin{quote}
The low-lying villages went under in the surge.
\end{quote}

Past (he/she …): \bn{সেই} \bn{নৌকা} \bn{তলিয়ে} \bn{গেল}, \emph{sei noukā tôliye gelô}.

\begin{cousinweb}
\textbf{\bn{তলিয়ে} \bn{যাওয়া}} (\emph{tôliye jāoyā}) is the -e form of \bn{তলানো}, to sink to the bottom, built on \bn{তল} — Sanskrit \emph{tala} (CDIAL 5731, which lists Bengali \emph{tal}, 'bottom, floor') — before \textbf{\bn{যাওয়া}}. \textbf{\bn{তলা}}, a storey, is from the same word.
\end{cousinweb}

\begin{grammarlens}[title={What you've built}]
One more word for telling what happened.
\end{grammarlens}

\subsection*{Your turn}
\begin{itemize}
\raggedright
  \item \emph{Say it:} \emph{tôliye jāoyā}
  \item \emph{Say it:} \emph{tôliye jāoyā}, once more
  \item \emph{Say it:} \emph{jôlochchhvāse nichu grāmgulo tôliye gelô.} — The low-lying villages went under in the surge.
  \item \emph{Recall:} say the Bengali for electricity, then the Bengali for a tidal surge, then say \emph{tôliye jāoyā} again
\end{itemize}

\begin{tcolorbox}[breakable,colback=teal!4,colframe=teal!35!black,title={Before you move on}]
What does \bn{তলিয়ে} \bn{যাওয়া} mean? (\textquotedblleft{}To go under; to be submerged\textquotedblright{}.) Say it once more.
\end{tcolorbox}
\section[\texorpdfstring{\bn{উঁচু} (ũchu)}{উঁচু (ũchu)}]{\bn{উঁচু} (ũchu) — high; raised}

\label{lesson:BN-C289-uchu}



\begin{quote}
Before the new one: say the Bengali for a tidal surge, then the Bengali for to go under.
\end{quote}

\subsection*{You'll want to know: \bn{উঁচু}}
\textbf{\bn{উঁচু}} — \emph{ũchu} — \textquotedblleft{}high; raised\textquotedblright{}.

\begin{quote}
\bn{সবাই} \bn{উঁচু} \bn{জায়গায়} \bn{আশ্রয়} \bn{নিল।}
\end{quote}

\begin{quote}
\emph{sôbāi ũchu jāygāy āshrôy nilô.}
\end{quote}

\begin{quote}
Everyone took shelter on higher ground.
\end{quote}

\begin{cousinweb}
\textbf{\bn{উঁচু}} (\emph{ũchu}) comes down from Sanskrit \emph{ucca}, 'tall, lofty' (CDIAL 1634, which lists Bengali \emph{uc} and \emph{ũc}); its opposite is \textbf{\bn{নিচু}}, low.
\end{cousinweb}

\begin{grammarlens}[title={What you've built}]
One more word for telling what happened.
\end{grammarlens}

\subsection*{Your turn}
\begin{itemize}
\raggedright
  \item \emph{Say it:} \emph{ũchu}
  \item \emph{Say it:} \emph{ũchu}, once more
  \item \emph{Say it:} \emph{sôbāi ũchu jāygāy āshrôy nilô.} — Everyone took shelter on higher ground.
  \item \emph{Recall:} say the Bengali for a tidal surge, then the Bengali for to go under, then say \emph{ũchu} again
\end{itemize}

\begin{tcolorbox}[breakable,colback=teal!4,colframe=teal!35!black,title={Before you move on}]
What does \bn{উঁচু} mean? (\textquotedblleft{}High; raised\textquotedblright{}.) Say it once more.
\end{tcolorbox}
\section[\texorpdfstring{\bn{ভেলা} (bhelā)}{ভেলা (bhelā)}]{\bn{ভেলা} (bhelā) — a raft}

\label{lesson:BN-C289-bhela}



\begin{quote}
Before the new one: say the Bengali for to go under, then the Bengali for high.
\end{quote}

\subsection*{You'll want to know: \bn{ভেলা}}
\textbf{\bn{ভেলা}} — \emph{bhelā} — \textquotedblleft{}a raft\textquotedblright{}.

\begin{quote}
\bn{জেলেরা} \bn{কলার} \bn{গাছ} \bn{দিয়ে} \bn{একটা} \bn{ভেলা} \bn{বানাল।}
\end{quote}

\begin{quote}
\emph{jelerā kôlār gāchh diye ekṭā bhelā bānālô.}
\end{quote}

\begin{quote}
The fishermen made a raft out of banana trees.
\end{quote}

\begin{cousinweb}
\textbf{\bn{ভেলা}} (\emph{bhelā}) comes down from Sanskrit \emph{bēḍā}, 'a boat', and its variant \emph{bhēla} (CDIAL 9308, which lists Bengali \emph{bhelā}, 'a raft').
\end{cousinweb}

\begin{grammarlens}[title={What you've built}]
One more word for telling what happened.
\end{grammarlens}

\subsection*{Your turn}
\begin{itemize}
\raggedright
  \item \emph{Say it:} \emph{bhelā}
  \item \emph{Say it:} \emph{bhelā}, once more
  \item \emph{Say it:} \emph{jelerā kôlār gāchh diye ekṭā bhelā bānālô.} — The fishermen made a raft out of banana trees.
  \item \emph{Recall:} say the Bengali for to go under, then the Bengali for high, then say \emph{bhelā} again
\end{itemize}

\begin{tcolorbox}[breakable,colback=teal!4,colframe=teal!35!black,title={Before you move on}]
What does \bn{ভেলা} mean? (\textquotedblleft{}A raft\textquotedblright{}.) Say it once more.
\end{tcolorbox}
\section[\texorpdfstring{\bn{প্লাবিত} \bn{হওয়া} (plābitô hôoyā)}{প্লাবিত হওয়া (plābitô hôoyā)}]{\bn{প্লাবিত} \bn{হওয়া} (plābitô hôoyā) — to be flooded; to be inundated}

\label{lesson:BN-C289-plabito-hooya}



\begin{quote}
Before the new one: say the Bengali for high, then the Bengali for a raft.
\end{quote}

\subsection*{You'll want to know: \bn{প্লাবিত} \bn{হওয়া}}
\textbf{\bn{প্লাবিত} \bn{হওয়া}} — \emph{plābitô hôoyā} — \textquotedblleft{}to be flooded; to be inundated\textquotedblright{}.

\begin{quote}
\bn{বন্যায়} \bn{উপকূলের} \bn{অনেক} \bn{গ্রাম} \bn{প্লাবিত} \bn{হল।}
\end{quote}

\begin{quote}
\emph{bônyāy upôkūler ônek grām plābitô hôlô.}
\end{quote}

\begin{quote}
Many villages on the coast were flooded.
\end{quote}

Past (he/she …): \bn{সেই} \bn{শহর} \bn{প্লাবিত} \bn{হল}, \emph{sei shôhôr plābitô hôlô}.

\begin{cousinweb}
\textbf{\bn{প্লাবিত} \bn{হওয়া}} (\emph{plābitô hôoyā}) puts Sanskrit \emph{plāvita}, 'made to overflow, deluged' (Monier-Williams), before \textbf{\bn{হওয়া}}: to be under the flood. It is a word of the news; talk says \bn{জলে} \bn{ডুবে} \bn{গেল}.
\end{cousinweb}

\begin{grammarlens}[title={What you've built}]
One more word for telling what happened.
\end{grammarlens}

\subsection*{Your turn}
\begin{itemize}
\raggedright
  \item \emph{Say it:} \emph{plābitô hôoyā}
  \item \emph{Say it:} \emph{plābitô hôoyā}, once more
  \item \emph{Say it:} \emph{bônyāy upôkūler ônek grām plābitô hôlô.} — Many villages on the coast were flooded.
  \item \emph{Recall:} say the Bengali for high, then the Bengali for a raft, then say \emph{plābitô hôoyā} again
\end{itemize}

\begin{tcolorbox}[breakable,colback=teal!4,colframe=teal!35!black,title={Before you move on}]
What does \bn{প্লাবিত} \bn{হওয়া} mean? (\textquotedblleft{}To be flooded; to be inundated\textquotedblright{}.) Say it once more.
\end{tcolorbox}
